
\documentclass{article}
\usepackage{colortbl}
\usepackage{makecell}
\usepackage{multirow}
\usepackage{supertabular}

\begin{document}

\newcounter{utterance}

\centering \large Interaction Transcript for game `clean\_up', experiment `2\_hard\_3obj\_el', episode 4 with deepseek{-}v4{-}pro.
\vspace{24pt}

{ \footnotesize  \setcounter{utterance}{1}
\setlength{\tabcolsep}{0pt}
\begin{supertabular}{c@{$\;$}|p{.15\linewidth}@{}p{.15\linewidth}p{.15\linewidth}p{.15\linewidth}p{.15\linewidth}p{.15\linewidth}}
    \# & $\;$A & \multicolumn{4}{c}{Game Master} & $\;\:$B\\
    \hline

    \theutterance \stepcounter{utterance}  
    & & \multicolumn{4}{p{0.6\linewidth}}{
        \cellcolor[rgb]{0.9,0.9,0.9}{
            \makecell[{{p{\linewidth}}}]{
                \texttt{\tiny{[P1$\langle$GM]}}
                \texttt{Είμαι ο αρχηγός του παιχνιδιού σας και παίζετε ένα συνεργατικό παιχνίδι με το ακόλουθο πλέγμα ως πίνακα παιχνιδιού:} \\
\\ 
\\ 
\texttt{ 1234567} \\
\texttt{╔══╤═╤══╗} \\
\texttt{║◌◌│◌│◌◌║ 1} \\
\texttt{║◌◌└─┼──╢ 2} \\
\texttt{║◌◌◌◌│P◌║ 3} \\
\texttt{║◌┌─┬┴┬─╢ 4} \\
\texttt{╟─┤◌│◌│◌║ 5} \\
\texttt{║◌└┬┴┬┘L║ 6} \\
\texttt{║◌◌│C│◌◌║ 7} \\
\texttt{╚══╧═╧══╝} \\
\\ 
\\ 
\\ 
\texttt{* Το άνω άκρο εμφανίζει συντεταγμένες x που αυξάνονται προς τα δεξιά και το δεξιό άκρο συντεταγμένες y που αυξάνονται προς τα κάτω.} \\
\texttt{* Τα ακόλουθα αντικείμενα τοποθετούνται τυχαία στο πλέγμα σας: 'C', 'L', 'P'.} \\
\\ 
\\ 
\texttt{Ο άλλος παίκτης βλέπει μια παραλλαγή του σκάφους του παιχνιδιού, όπου τα αντικείμενα τοποθετούνται σε διαφορετικές τυχαίες θέσεις. Δεν μπορείτε να δείτε το ταμπλό του άλλου παίκτη και εκείνος δεν μπορεί να δει το δικό σας.} \\
\\ 
\\ 
\texttt{**Στόχος:**} \\
\\ 
\texttt{Και οι δύο παίκτες πρέπει να μετακινήσουν τα αντικείμενα στο αντίστοιχο φόντο τους, έτσι ώστε πανομοιότυπα αντικείμενα να καταλήξουν στις ίδιες συντεταγμένες. Πρέπει να επικοινωνήσετε με τον άλλο παίκτη για να συμφωνήσετε σε μια κοινή διαμόρφωση του στόχου.} \\
\\ 
\\ 
\texttt{**Κανόνες:**} \\
\\ 
\texttt{* Σε κάθε γύρο, μπορείτε να στείλετε ακριβώς μία από τις ακόλουθες δύο εντολές:} \\
\texttt{ 1. `ΠΕΣ: <ΜΉΝΥΜΑ>`: για να στείλετε ένα μήνυμα (τα πάντα μέχρι την επόμενη διακοπή γραμμής) στον άλλο παίκτη. Θα το προωθήσω στον παρτενέρ σας.} \\
\texttt{2. `ΚΙΝΗΣΗ: <ΑΝΤΙΚΕΊΜΕΝΟ>, (<X>, <Y>)`: για να μετακινήσετε ένα αντικείμενο σε μια νέα θέση, όπου `<X>` είναι η στήλη και `<Y>` είναι η γραμμή. Θα σας ενημερώσω αν η μετακίνησή σας ήταν επιτυχής ή όχι.} \\
\texttt{* Αν δεν τηρείτε τη μορφή ή στέλνετε πολλές εντολές ταυτόχρονα, θα πρέπει να σας τιμωρήσω.} \\
\texttt{* Εάν και οι δύο παίκτες συγκεντρώσουν περισσότερες από 8 ποινές, χάνετε και οι δύο το παιχνίδι.} \\
\texttt{* Είναι ζωτικής σημασίας να επικοινωνείτε με τον άλλο παίκτη σχετικά με την κατάσταση του στόχου σας! Ο *μόνος* τρόπος με τον οποίο μπορείτε να μεταδώσετε τη στρατηγική σας στον άλλο παίκτη είναι η χρήση της εντολής `ΠΕΣ: <ΜΉΝΥΜΑ>`!} \\
\\ 
\\ 
\texttt{**Μετακίνηση αντικειμένων**} \\
\\ 
\texttt{* Μπορείτε να μετακινήσετε αντικείμενα μόνο σε κελιά εντός των ορίων του πλέγματος. Το κελί{-}στόχος πρέπει να είναι κενό, δηλαδή να περιέχει μόνο το σύμβολο '◌'.} \\
\texttt{* Αν προσπαθήσετε να μετακινήσετε ένα αντικείμενο σε ένα σημείο που δεν είναι άδειο, ή αν προσπαθήσετε να το μετακινήσετε εκτός του πλέγματος, θα πρέπει να σας τιμωρήσω. Μπορείτε να προσπαθήσετε ξανά.} \\
\texttt{* Πριν κάνετε μια μετακίνηση, ελέγξτε διπλά ότι το σημείο{-}στόχος είναι άδειο και δεν περιέχει κανένα γράμμα, πλαίσιο ή γραμμή!} \\
\\ 
\\ 
\texttt{**Τέλος του παιχνιδιού**} \\
\\ 
\texttt{Αν νομίζετε ότι έχετε πετύχει το στόχο της ευθυγράμμισης όλων των αντικειμένων, μπορείτε να ζητήσετε από τον άλλο παίκτη να τερματίσει το παιχνίδι στέλνοντας `ΠΕΣ: τελείωσες;`. Αν ο άλλος παίκτης σας ζητήσει να τερματίσετε το παιχνίδι και απαντήσετε με `ΠΕΣ: τελείωσα!`, το παιχνίδι θα τελειώσει.} \\
\\ 
\texttt{Και οι δύο παίκτες κερδίζουν αν το παιχνίδι τελειώσει μέσα σε 12 γύρους, όπου ένας γύρος ορίζεται ως δύο παίκτες που στέλνουν από μία έγκυρη εντολή.} \\
\\ 
\texttt{**Βαθμολογία:**} \\
\\ 
\texttt{Όσο πιο κοντά βρίσκονται τα ίδια αντικείμενα και στους δύο πίνακες του παιχνιδιού, τόσο περισσότερους πόντους παίρνετε. Οι ποινές μειώνουν τους πόντους σας. Μπορείτε να καταρρίψετε το ρεκόρ;} \\
\\ 
\\ 
\texttt{Παρακαλώ στείλτε ένα μήνυμα στον άλλο παίκτη για να ξεκινήσετε το παιχνίδι!} \\
            }
        }
    }
    & \\ \\

    \theutterance \stepcounter{utterance}  
    & \multicolumn{4}{p{0.6\linewidth}}{
        \cellcolor[rgb]{0.9,0.9,0.9}{
            \makecell[{{p{\linewidth}}}]{
                \texttt{\tiny{[P1$\rangle$GM]}}
                \texttt{ΠΕΣ: Γεια! Ας συντονιστούμε για να βρούμε κοινή θέση για τα C, L, P. Στον πίνακά μου βλέπω: C στη θέση (4,7), L στη θέση (7,6), P στη θέση (6,3). Εσύ τι θέσεις έχεις;} \\
            }
        }
    }
    & & \\ \\

    \theutterance \stepcounter{utterance}  
    & & & \multicolumn{2}{p{0.3\linewidth}}{
        \cellcolor[rgb]{0.9,0.9,0.9}{
            \makecell[{{p{\linewidth}}}]{
                \texttt{\tiny{[GM$|$GM]}}
                \texttt{ΠΕΣ: Γεια! Ας συντονιστούμε για να βρούμε κοινή θέση για τα C, L, P. Στον πίνακά μου βλέπω: C στη θέση (4,7), L στη θέση (7,6), P στη θέση (6,3). Εσύ τι θέσεις έχεις;} \\
            }
        }
    }
    & & \\ \\

    \theutterance \stepcounter{utterance}  
    & & \multicolumn{4}{p{0.6\linewidth}}{
        \cellcolor[rgb]{0.9,0.9,0.9}{
            \makecell[{{p{\linewidth}}}]{
                \texttt{\tiny{[GM$\rangle$P2]}}
                \texttt{Είμαι ο αρχηγός του παιχνιδιού σας και παίζετε ένα συνεργατικό παιχνίδι με το ακόλουθο πλέγμα ως πίνακα παιχνιδιού:} \\
\\ 
\\ 
\texttt{ 1234567} \\
\texttt{╔══╤═╤══╗} \\
\texttt{║◌P│◌│◌◌║ 1} \\
\texttt{║◌L└─┼──╢ 2} \\
\texttt{║◌◌◌◌│◌◌║ 3} \\
\texttt{║◌┌─┬┴┬─╢ 4} \\
\texttt{╟─┤◌│◌│◌║ 5} \\
\texttt{║◌└┬┴┬┘◌║ 6} \\
\texttt{║◌◌│◌│◌C║ 7} \\
\texttt{╚══╧═╧══╝} \\
\\ 
\\ 
\\ 
\texttt{* Το άνω άκρο εμφανίζει συντεταγμένες x που αυξάνονται προς τα δεξιά και το δεξιό άκρο συντεταγμένες y που αυξάνονται προς τα κάτω.} \\
\texttt{* Τα ακόλουθα αντικείμενα τοποθετούνται τυχαία στο πλέγμα σας: 'C', 'L', 'P'.} \\
\\ 
\\ 
\texttt{Ο άλλος παίκτης βλέπει μια παραλλαγή του σκάφους του παιχνιδιού, όπου τα αντικείμενα τοποθετούνται σε διαφορετικές τυχαίες θέσεις. Δεν μπορείτε να δείτε το ταμπλό του άλλου παίκτη και εκείνος δεν μπορεί να δει το δικό σας.} \\
\\ 
\\ 
\texttt{**Στόχος:**} \\
\\ 
\texttt{Και οι δύο παίκτες πρέπει να μετακινήσουν τα αντικείμενα στο αντίστοιχο φόντο τους, έτσι ώστε πανομοιότυπα αντικείμενα να καταλήξουν στις ίδιες συντεταγμένες. Πρέπει να επικοινωνήσετε με τον άλλο παίκτη για να συμφωνήσετε σε μια κοινή διαμόρφωση του στόχου.} \\
\\ 
\\ 
\texttt{**Κανόνες:**} \\
\\ 
\texttt{* Σε κάθε γύρο, μπορείτε να στείλετε ακριβώς μία από τις ακόλουθες δύο εντολές:} \\
\texttt{ 1. `ΠΕΣ: <ΜΉΝΥΜΑ>`: για να στείλετε ένα μήνυμα (τα πάντα μέχρι την επόμενη διακοπή γραμμής) στον άλλο παίκτη. Θα το προωθήσω στον παρτενέρ σας.} \\
\texttt{2. `ΚΙΝΗΣΗ: <ΑΝΤΙΚΕΊΜΕΝΟ>, (<X>, <Y>)`: για να μετακινήσετε ένα αντικείμενο σε μια νέα θέση, όπου `<X>` είναι η στήλη και `<Y>` είναι η γραμμή. Θα σας ενημερώσω αν η μετακίνησή σας ήταν επιτυχής ή όχι.} \\
\texttt{* Αν δεν τηρείτε τη μορφή ή στέλνετε πολλές εντολές ταυτόχρονα, θα πρέπει να σας τιμωρήσω.} \\
\texttt{* Εάν και οι δύο παίκτες συγκεντρώσουν περισσότερες από 8 ποινές, χάνετε και οι δύο το παιχνίδι.} \\
\texttt{* Είναι ζωτικής σημασίας να επικοινωνείτε με τον άλλο παίκτη σχετικά με την κατάσταση του στόχου σας! Ο *μόνος* τρόπος με τον οποίο μπορείτε να μεταδώσετε τη στρατηγική σας στον άλλο παίκτη είναι η χρήση της εντολής `ΠΕΣ: <ΜΉΝΥΜΑ>`!} \\
\\ 
\\ 
\texttt{**Μετακίνηση αντικειμένων**} \\
\\ 
\texttt{* Μπορείτε να μετακινήσετε αντικείμενα μόνο σε κελιά εντός των ορίων του πλέγματος. Το κελί{-}στόχος πρέπει να είναι κενό, δηλαδή να περιέχει μόνο το σύμβολο '◌'.} \\
\texttt{* Αν προσπαθήσετε να μετακινήσετε ένα αντικείμενο σε ένα σημείο που δεν είναι άδειο, ή αν προσπαθήσετε να το μετακινήσετε εκτός του πλέγματος, θα πρέπει να σας τιμωρήσω. Μπορείτε να προσπαθήσετε ξανά.} \\
\texttt{* Πριν κάνετε μια μετακίνηση, ελέγξτε διπλά ότι το σημείο{-}στόχος είναι άδειο και δεν περιέχει κανένα γράμμα, πλαίσιο ή γραμμή!} \\
\\ 
\\ 
\texttt{**Τέλος του παιχνιδιού**} \\
\\ 
\texttt{Αν νομίζετε ότι έχετε πετύχει το στόχο της ευθυγράμμισης όλων των αντικειμένων, μπορείτε να ζητήσετε από τον άλλο παίκτη να τερματίσει το παιχνίδι στέλνοντας `ΠΕΣ: τελείωσες;`. Αν ο άλλος παίκτης σας ζητήσει να τερματίσετε το παιχνίδι και απαντήσετε με `ΠΕΣ: τελείωσα!`, το παιχνίδι θα τελειώσει.} \\
\\ 
\texttt{Και οι δύο παίκτες κερδίζουν αν το παιχνίδι τελειώσει μέσα σε 12 γύρους, όπου ένας γύρος ορίζεται ως δύο παίκτες που στέλνουν από μία έγκυρη εντολή.} \\
\\ 
\texttt{**Βαθμολογία:**} \\
\\ 
\texttt{Όσο πιο κοντά βρίσκονται τα ίδια αντικείμενα και στους δύο πίνακες του παιχνιδιού, τόσο περισσότερους πόντους παίρνετε. Οι ποινές μειώνουν τους πόντους σας. Μπορείτε να καταρρίψετε το ρεκόρ;} \\
\\ 
\texttt{Ο άλλος παίκτης ξεκίνησε το παιχνίδι στέλνοντας αυτό το μήνυμα:} \\
\\ 
\texttt{"Γεια! Ας συντονιστούμε για να βρούμε κοινή θέση για τα C, L, P. Στον πίνακά μου βλέπω: C στη θέση (4,7), L στη θέση (7,6), P στη θέση (6,3). Εσύ τι θέσεις έχεις;"} \\
\\ 
\texttt{Ποια είναι η πρώτη σας εντολή;} \\
            }
        }
    }
    & \\ \\

    \theutterance \stepcounter{utterance}  
    & & & \multicolumn{4}{p{0.6\linewidth}}{
        \cellcolor[rgb]{0.9,0.9,0.9}{
            \makecell[{{p{\linewidth}}}]{
                \texttt{\tiny{[GM$\langle$P2]}}
                \texttt{Σκέφτομαι βήμα προς βήμα:} \\
\\ 
\texttt{1. Πρώτα, πρέπει να καταλάβω τις θέσεις των αντικειμένων στον δικό μου πίνακα. Από το πλέγμα που μου δόθηκε:} \\
\\ 
\texttt{Διαβάζοντας το πλέγμα (x: 1{-}7 από αριστερά προς τα δεξιά, y: 1{-}7 από πάνω προς τα κάτω):} \\
\\ 
\texttt{{-} Σειρά 1: ◌P│◌│◌◌ → P στη θέση (2,1)} \\
\texttt{{-} Σειρά 2: ◌L└─┼── → L στη θέση (2,2)} \\
\texttt{{-} Σειρά 3: ◌◌◌◌│◌◌} \\
\texttt{{-} Σειρά 4: ◌┌─┬┴┬─} \\
\texttt{{-} Σειρά 5: ─┤◌│◌│◌} \\
\texttt{{-} Σειρά 6: ◌└┬┴┬┘◌} \\
\texttt{{-} Σειρά 7: ◌◌│◌│◌C → C στη θέση (7,7)} \\
\\ 
\texttt{2. Οι θέσεις μου είναι:} \\
\texttt{   {-} C: (7,7)} \\
\texttt{   {-} L: (2,2)} \\
\texttt{   {-} P: (2,1)} \\
\\ 
\texttt{3. Οι θέσεις του άλλου παίκτη είναι:} \\
\texttt{   {-} C: (4,7)} \\
\texttt{   {-} L: (7,6)} \\
\texttt{   {-} P: (6,3)} \\
\\ 
\texttt{4. Πρέπει να επικοινωνήσω τις θέσεις μου και να συμφωνήσουμε σε μια κοινή διαμόρφωση. Θα στείλω πρώτα τα δεδομένα μου.} \\
\\ 
\texttt{5. Η πρώτη μου εντολή θα είναι ένα μήνυμα που θα αναφέρει τις θέσεις μου.} \\
\\ 
\texttt{ΠΕΣ: Γεια! Στον δικό μου πίνακα: C στη θέση (7,7), L στη θέση (2,2), P στη θέση (2,1). Πρέπει να συμφωνήσουμε σε κοινές θέσεις{-}στόχους. Προτείνω να βρούμε άδεια κελιά που να είναι προσβάσιμα και για τους δύο. Έχεις προτάσεις;} \\
            }
        }
    }
     \\ \\

    \theutterance \stepcounter{utterance}  
    & & & \multicolumn{2}{p{0.3\linewidth}}{
        \cellcolor[rgb]{0.9,0.9,0.9}{
            \makecell[{{p{\linewidth}}}]{
                \texttt{\tiny{[GM$|$GM]}}
                \texttt{Σκέφτομαι βήμα προς βήμα:} \\
\\ 
\texttt{1. Πρώτα, πρέπει να καταλάβω τις θέσεις των αντικειμένων στον δικό μου πίνακα. Από το πλέγμα που μου δόθηκε:} \\
\\ 
\texttt{Διαβάζοντας το πλέγμα (x: 1{-}7 από αριστερά προς τα δεξιά, y: 1{-}7 από πάνω προς τα κάτω):} \\
\\ 
\texttt{{-} Σειρά 1: ◌P│◌│◌◌ → P στη θέση (2,1)} \\
\texttt{{-} Σειρά 2: ◌L└─┼── → L στη θέση (2,2)} \\
\texttt{{-} Σειρά 3: ◌◌◌◌│◌◌} \\
\texttt{{-} Σειρά 4: ◌┌─┬┴┬─} \\
\texttt{{-} Σειρά 5: ─┤◌│◌│◌} \\
\texttt{{-} Σειρά 6: ◌└┬┴┬┘◌} \\
\texttt{{-} Σειρά 7: ◌◌│◌│◌C → C στη θέση (7,7)} \\
\\ 
\texttt{2. Οι θέσεις μου είναι:} \\
\texttt{   {-} C: (7,7)} \\
\texttt{   {-} L: (2,2)} \\
\texttt{   {-} P: (2,1)} \\
\\ 
\texttt{3. Οι θέσεις του άλλου παίκτη είναι:} \\
\texttt{   {-} C: (4,7)} \\
\texttt{   {-} L: (7,6)} \\
\texttt{   {-} P: (6,3)} \\
\\ 
\texttt{4. Πρέπει να επικοινωνήσω τις θέσεις μου και να συμφωνήσουμε σε μια κοινή διαμόρφωση. Θα στείλω πρώτα τα δεδομένα μου.} \\
\\ 
\texttt{5. Η πρώτη μου εντολή θα είναι ένα μήνυμα που θα αναφέρει τις θέσεις μου.} \\
\\ 
\texttt{ΠΕΣ: Γεια! Στον δικό μου πίνακα: C στη θέση (7,7), L στη θέση (2,2), P στη θέση (2,1). Πρέπει να συμφωνήσουμε σε κοινές θέσεις{-}στόχους. Προτείνω να βρούμε άδεια κελιά που να είναι προσβάσιμα και για τους δύο. Έχεις προτάσεις;} \\
            }
        }
    }
    & & \\ \\

    \theutterance \stepcounter{utterance}  
    & & & \multicolumn{2}{p{0.3\linewidth}}{
        \cellcolor[rgb]{0.9,0.9,0.9}{
            \makecell[{{p{\linewidth}}}]{
                \texttt{\tiny{[GM$|$GM]}}
                \texttt{Invalid format: head is not empty: 'Σκέφτομαι βήμα προς βήμα:} \\
\\ 
\texttt{1. Πρώτα, πρέπει να καταλάβω τις θέσεις των αντικειμένων στον δικό μου πίνακα. Από το πλέγμα που μου δόθηκε:} \\
\\ 
\texttt{Διαβάζοντας το πλέγμα (x: 1{-}7 από αριστερά προς τα δεξιά, y: 1{-}7 από πάνω προς τα κάτω):} \\
\\ 
\texttt{{-} Σειρά 1: ◌P│◌│◌◌ → P στη θέση (2,1)} \\
\texttt{{-} Σειρά 2: ◌L└─┼── → L στη θέση (2,2)} \\
\texttt{{-} Σειρά 3: ◌◌◌◌│◌◌} \\
\texttt{{-} Σειρά 4: ◌┌─┬┴┬─} \\
\texttt{{-} Σειρά 5: ─┤◌│◌│◌} \\
\texttt{{-} Σειρά 6: ◌└┬┴┬┘◌} \\
\texttt{{-} Σειρά 7: ◌◌│◌│◌C → C στη θέση (7,7)} \\
\\ 
\texttt{2. Οι θέσεις μου είναι:} \\
\texttt{   {-} C: (7,7)} \\
\texttt{   {-} L: (2,2)} \\
\texttt{   {-} P: (2,1)} \\
\\ 
\texttt{3. Οι θέσεις του άλλου παίκτη είναι:} \\
\texttt{   {-} C: (4,7)} \\
\texttt{   {-} L: (7,6)} \\
\texttt{   {-} P: (6,3)} \\
\\ 
\texttt{4. Πρέπει να επικοινωνήσω τις θέσεις μου και να συμφωνήσουμε σε μια κοινή διαμόρφωση. Θα στείλω πρώτα τα δεδομένα μου.} \\
\\ 
\texttt{5. Η πρώτη μου εντολή θα είναι ένα μήνυμα που θα αναφέρει τις θέσεις μου.} \\
\\ 
\texttt{'} \\
            }
        }
    }
    & & \\ \\

    \theutterance \stepcounter{utterance}  
    & & \multicolumn{4}{p{0.6\linewidth}}{
        \cellcolor[rgb]{0.9,0.9,0.9}{
            \makecell[{{p{\linewidth}}}]{
                \texttt{\tiny{[GM$\rangle$P2]}}
                \texttt{Ποινή: Το μήνυμά σας δεν πρέπει να περιέχει τίποτα πριν από την εντολή!} \\
\\ 
\texttt{Βεβαιωθείτε ότι η απάντησή σας περιέχει μόνο `ΠΕΣ: <ΜΉΝΥΜΑ>` ή `ΚΙΝΗΣΗ: <ΑΝΤΙΚΕΊΜΕΝΟ>, (<X>, <Y>)`, και τίποτα άλλο!} \\
\\ 
\texttt{Έχετε συσσωρεύσει συλλογικά 1 από 8 ποινές.} \\
\texttt{Παρακαλώ δοκιμάστε ξανά!} \\
            }
        }
    }
    & \\ \\

    \theutterance \stepcounter{utterance}  
    & & & \multicolumn{4}{p{0.6\linewidth}}{
        \cellcolor[rgb]{0.9,0.9,0.9}{
            \makecell[{{p{\linewidth}}}]{
                \texttt{\tiny{[GM$\langle$P2]}}
                \texttt{ΠΕΣ: Γεια! Στον δικό μου πίνακα: C στη θέση (7,7), L στη θέση (2,2), P στη θέση (2,1). Πρέπει να συμφωνήσουμε σε κοινές θέσεις{-}στόχους. Προτείνω να βρούμε άδεια κελιά που να είναι προσβάσιμα και για τους δύο. Έχεις προτάσεις;} \\
            }
        }
    }
     \\ \\

    \theutterance \stepcounter{utterance}  
    & & & \multicolumn{2}{p{0.3\linewidth}}{
        \cellcolor[rgb]{0.9,0.9,0.9}{
            \makecell[{{p{\linewidth}}}]{
                \texttt{\tiny{[GM$|$GM]}}
                \texttt{ΠΕΣ: Γεια! Στον δικό μου πίνακα: C στη θέση (7,7), L στη θέση (2,2), P στη θέση (2,1). Πρέπει να συμφωνήσουμε σε κοινές θέσεις{-}στόχους. Προτείνω να βρούμε άδεια κελιά που να είναι προσβάσιμα και για τους δύο. Έχεις προτάσεις;} \\
            }
        }
    }
    & & \\ \\

    \theutterance \stepcounter{utterance}  
    & & \multicolumn{4}{p{0.6\linewidth}}{
        \cellcolor[rgb]{0.9,0.9,0.9}{
            \makecell[{{p{\linewidth}}}]{
                \texttt{\tiny{[P1$\langle$GM]}}
                \texttt{Το μήνυμά σας μεταφέρθηκε στον άλλο παίκτη.} \\
\\ 
\\ 
\texttt{Αυτή τη στιγμή παίζετε τον γύρο 1 από συνολικά 12.} \\
\texttt{Έχετε συσσωρεύσει συλλογικά 1 από 8 ποινές.} \\
\texttt{Ο άλλος παίκτης έστειλε αυτό το μήνυμα:} \\
\\ 
\\ 
\texttt{"Γεια! Στον δικό μου πίνακα: C στη θέση (7,7), L στη θέση (2,2), P στη θέση (2,1). Πρέπει να συμφωνήσουμε σε κοινές θέσεις{-}στόχους. Προτείνω να βρούμε άδεια κελιά που να είναι προσβάσιμα και για τους δύο. Έχεις προτάσεις;"} \\
\\ 
\\ 
\texttt{Ποια είναι η επόμενη εντολή σας;} \\
            }
        }
    }
    & \\ \\

    \theutterance \stepcounter{utterance}  
    & \multicolumn{4}{p{0.6\linewidth}}{
        \cellcolor[rgb]{0.9,0.9,0.9}{
            \makecell[{{p{\linewidth}}}]{
                \texttt{\tiny{[P1$\rangle$GM]}}
                \texttt{ΠΕΣ: Ας προτείνουμε συγκεκριμένες θέσεις{-}στόχους. Βλέπω αρκετά άδεια κελιά. Τι θα έλεγες για C στο (3,3), L στο (5,5), P στο (4,2); Ελέγχω αν αυτά είναι άδεια στον πίνακά μου... Το (3,3) είναι άδειο, το (5,5) έχει ένα πλαίσιο '┤' οπότε δεν είναι άδειο, το (4,2) είναι άδειο. Χρειαζόμαστε διόρθωση. Εσύ τι κενά βλέπεις;} \\
            }
        }
    }
    & & \\ \\

    \theutterance \stepcounter{utterance}  
    & & & \multicolumn{2}{p{0.3\linewidth}}{
        \cellcolor[rgb]{0.9,0.9,0.9}{
            \makecell[{{p{\linewidth}}}]{
                \texttt{\tiny{[GM$|$GM]}}
                \texttt{ΠΕΣ: Ας προτείνουμε συγκεκριμένες θέσεις{-}στόχους. Βλέπω αρκετά άδεια κελιά. Τι θα έλεγες για C στο (3,3), L στο (5,5), P στο (4,2); Ελέγχω αν αυτά είναι άδεια στον πίνακά μου... Το (3,3) είναι άδειο, το (5,5) έχει ένα πλαίσιο '┤' οπότε δεν είναι άδειο, το (4,2) είναι άδειο. Χρειαζόμαστε διόρθωση. Εσύ τι κενά βλέπεις;} \\
            }
        }
    }
    & & \\ \\

    \theutterance \stepcounter{utterance}  
    & & \multicolumn{4}{p{0.6\linewidth}}{
        \cellcolor[rgb]{0.9,0.9,0.9}{
            \makecell[{{p{\linewidth}}}]{
                \texttt{\tiny{[GM$\rangle$P2]}}
                \texttt{Το μήνυμά σας μεταφέρθηκε στον άλλο παίκτη.} \\
\\ 
\\ 
\texttt{Αυτή τη στιγμή παίζετε τον γύρο 2 από συνολικά 12.} \\
\texttt{Έχετε συσσωρεύσει συλλογικά 1 από 8 ποινές.} \\
\texttt{Ο άλλος παίκτης έστειλε αυτό το μήνυμα:} \\
\\ 
\\ 
\texttt{"Ας προτείνουμε συγκεκριμένες θέσεις{-}στόχους. Βλέπω αρκετά άδεια κελιά. Τι θα έλεγες για C στο (3,3), L στο (5,5), P στο (4,2); Ελέγχω αν αυτά είναι άδεια στον πίνακά μου... Το (3,3) είναι άδειο, το (5,5) έχει ένα πλαίσιο '┤' οπότε δεν είναι άδειο, το (4,2) είναι άδειο. Χρειαζόμαστε διόρθωση. Εσύ τι κενά βλέπεις;"} \\
\\ 
\\ 
\texttt{Ποια είναι η επόμενη εντολή σας;} \\
            }
        }
    }
    & \\ \\

    \theutterance \stepcounter{utterance}  
    & & & \multicolumn{4}{p{0.6\linewidth}}{
        \cellcolor[rgb]{0.9,0.9,0.9}{
            \makecell[{{p{\linewidth}}}]{
                \texttt{\tiny{[GM$\langle$P2]}}
                \texttt{ΠΕΣ: Στον πίνακά μου, ελέγχω τα κελιά: (3,3) είναι άδειο, (5,5) περιέχει "┴" οπότε δεν είναι άδειο, (4,2) περιέχει "└" οπότε δεν είναι άδειο. Χρειαζόμαστε στόχους που να είναι άδειοι και στους δύο πίνακες. Ας βρούμε τομές από άδεια κελιά. Μπορώ να δω τα δικά μου άδεια σημεία: (1,1), (1,2), (1,3), (1,4), (1,6), (1,7), (3,1), (4,1), (5,1), (6,1), (7,1), (3,2), (4,3), (5,3), (6,3), (7,3), (2,4), (5,4), (7,4), (2,5), (4,5), (6,5), (7,5), (1,6), (2,6), (7,6), (2,7), (3,7), (5,7), (6,7). Πες μου τα δικά σου άδεια κελιά να βρούμε κοινά!} \\
            }
        }
    }
     \\ \\

    \theutterance \stepcounter{utterance}  
    & & & \multicolumn{2}{p{0.3\linewidth}}{
        \cellcolor[rgb]{0.9,0.9,0.9}{
            \makecell[{{p{\linewidth}}}]{
                \texttt{\tiny{[GM$|$GM]}}
                \texttt{ΠΕΣ: Στον πίνακά μου, ελέγχω τα κελιά: (3,3) είναι άδειο, (5,5) περιέχει "┴" οπότε δεν είναι άδειο, (4,2) περιέχει "└" οπότε δεν είναι άδειο. Χρειαζόμαστε στόχους που να είναι άδειοι και στους δύο πίνακες. Ας βρούμε τομές από άδεια κελιά. Μπορώ να δω τα δικά μου άδεια σημεία: (1,1), (1,2), (1,3), (1,4), (1,6), (1,7), (3,1), (4,1), (5,1), (6,1), (7,1), (3,2), (4,3), (5,3), (6,3), (7,3), (2,4), (5,4), (7,4), (2,5), (4,5), (6,5), (7,5), (1,6), (2,6), (7,6), (2,7), (3,7), (5,7), (6,7). Πες μου τα δικά σου άδεια κελιά να βρούμε κοινά!} \\
            }
        }
    }
    & & \\ \\

    \theutterance \stepcounter{utterance}  
    & & \multicolumn{4}{p{0.6\linewidth}}{
        \cellcolor[rgb]{0.9,0.9,0.9}{
            \makecell[{{p{\linewidth}}}]{
                \texttt{\tiny{[P1$\langle$GM]}}
                \texttt{Το μήνυμά σας μεταφέρθηκε στον άλλο παίκτη.} \\
\\ 
\\ 
\texttt{Αυτή τη στιγμή παίζετε τον γύρο 2 από συνολικά 12.} \\
\texttt{Έχετε συσσωρεύσει συλλογικά 1 από 8 ποινές.} \\
\texttt{Ο άλλος παίκτης έστειλε αυτό το μήνυμα:} \\
\\ 
\\ 
\texttt{"Στον πίνακά μου, ελέγχω τα κελιά: (3,3) είναι άδειο, (5,5) περιέχει "┴" οπότε δεν είναι άδειο, (4,2) περιέχει "└" οπότε δεν είναι άδειο. Χρειαζόμαστε στόχους που να είναι άδειοι και στους δύο πίνακες. Ας βρούμε τομές από άδεια κελιά. Μπορώ να δω τα δικά μου άδεια σημεία: (1,1), (1,2), (1,3), (1,4), (1,6), (1,7), (3,1), (4,1), (5,1), (6,1), (7,1), (3,2), (4,3), (5,3), (6,3), (7,3), (2,4), (5,4), (7,4), (2,5), (4,5), (6,5), (7,5), (1,6), (2,6), (7,6), (2,7), (3,7), (5,7), (6,7). Πες μου τα δικά σου άδεια κελιά να βρούμε κοινά!"} \\
\\ 
\\ 
\texttt{Ποια είναι η επόμενη εντολή σας;} \\
            }
        }
    }
    & \\ \\

    \theutterance \stepcounter{utterance}  
    & \multicolumn{4}{p{0.6\linewidth}}{
        \cellcolor[rgb]{0.9,0.9,0.9}{
            \makecell[{{p{\linewidth}}}]{
                \texttt{\tiny{[P1$\rangle$GM]}}
                \texttt{ΠΕΣ: Ωραία ιδέα! Εδώ είναι τα δικά μου άδεια κελιά: (1,1), (1,4), (1,5), (1,6), (1,7), (2,1), (2,4), (2,5), (2,7), (3,1), (3,2), (3,4), (3,5), (3,7), (4,1), (4,2), (4,3), (4,4), (4,5), (4,6), (5,1), (5,2), (5,3), (5,5), (5,7), (6,1), (6,4), (6,5), (6,6), (6,7), (7,1), (7,2), (7,3), (7,4), (7,5). Ας βρούμε την τομή! Κοινά άδεια φαίνονται: (1,1), (1,4), (1,6), (1,7), (2,4), (2,5), (2,7), (3,1), (3,7), (4,1), (4,3), (5,1), (5,3), (5,5), (6,1), (6,4), (6,5), (7,1), (7,3), (7,4), (7,5). Από αυτά, προτείνω: C στο (3,7), L στο (1,4), P στο (7,4). Ελέγχεις αν είναι άδεια στον δικό σου;} \\
            }
        }
    }
    & & \\ \\

    \theutterance \stepcounter{utterance}  
    & & & \multicolumn{2}{p{0.3\linewidth}}{
        \cellcolor[rgb]{0.9,0.9,0.9}{
            \makecell[{{p{\linewidth}}}]{
                \texttt{\tiny{[GM$|$GM]}}
                \texttt{ΠΕΣ: Ωραία ιδέα! Εδώ είναι τα δικά μου άδεια κελιά: (1,1), (1,4), (1,5), (1,6), (1,7), (2,1), (2,4), (2,5), (2,7), (3,1), (3,2), (3,4), (3,5), (3,7), (4,1), (4,2), (4,3), (4,4), (4,5), (4,6), (5,1), (5,2), (5,3), (5,5), (5,7), (6,1), (6,4), (6,5), (6,6), (6,7), (7,1), (7,2), (7,3), (7,4), (7,5). Ας βρούμε την τομή! Κοινά άδεια φαίνονται: (1,1), (1,4), (1,6), (1,7), (2,4), (2,5), (2,7), (3,1), (3,7), (4,1), (4,3), (5,1), (5,3), (5,5), (6,1), (6,4), (6,5), (7,1), (7,3), (7,4), (7,5). Από αυτά, προτείνω: C στο (3,7), L στο (1,4), P στο (7,4). Ελέγχεις αν είναι άδεια στον δικό σου;} \\
            }
        }
    }
    & & \\ \\

    \theutterance \stepcounter{utterance}  
    & & \multicolumn{4}{p{0.6\linewidth}}{
        \cellcolor[rgb]{0.9,0.9,0.9}{
            \makecell[{{p{\linewidth}}}]{
                \texttt{\tiny{[GM$\rangle$P2]}}
                \texttt{Το μήνυμά σας μεταφέρθηκε στον άλλο παίκτη.} \\
\\ 
\\ 
\texttt{Αυτή τη στιγμή παίζετε τον γύρο 3 από συνολικά 12.} \\
\texttt{Έχετε συσσωρεύσει συλλογικά 1 από 8 ποινές.} \\
\texttt{Ο άλλος παίκτης έστειλε αυτό το μήνυμα:} \\
\\ 
\\ 
\texttt{"Ωραία ιδέα! Εδώ είναι τα δικά μου άδεια κελιά: (1,1), (1,4), (1,5), (1,6), (1,7), (2,1), (2,4), (2,5), (2,7), (3,1), (3,2), (3,4), (3,5), (3,7), (4,1), (4,2), (4,3), (4,4), (4,5), (4,6), (5,1), (5,2), (5,3), (5,5), (5,7), (6,1), (6,4), (6,5), (6,6), (6,7), (7,1), (7,2), (7,3), (7,4), (7,5). Ας βρούμε την τομή! Κοινά άδεια φαίνονται: (1,1), (1,4), (1,6), (1,7), (2,4), (2,5), (2,7), (3,1), (3,7), (4,1), (4,3), (5,1), (5,3), (5,5), (6,1), (6,4), (6,5), (7,1), (7,3), (7,4), (7,5). Από αυτά, προτείνω: C στο (3,7), L στο (1,4), P στο (7,4). Ελέγχεις αν είναι άδεια στον δικό σου;"} \\
\\ 
\\ 
\texttt{Ποια είναι η επόμενη εντολή σας;} \\
            }
        }
    }
    & \\ \\

    \theutterance \stepcounter{utterance}  
    & & & \multicolumn{4}{p{0.6\linewidth}}{
        \cellcolor[rgb]{0.9,0.9,0.9}{
            \makecell[{{p{\linewidth}}}]{
                \texttt{\tiny{[GM$\langle$P2]}}
                \texttt{ΠΕΣ: Έλεγξα τις προτάσεις σου στον πίνακά μου: (3,7) είναι άδειο, (1,4) είναι άδειο, (7,4) είναι άδειο. Συμφωνώ με C στο (3,7), L στο (1,4), P στο (7,4). Να ξεκινήσω τις μετακινήσεις; Αν συμφωνείς, πες μου να ξεκινήσουμε και εγώ θα μετακινήσω πρώτα το P από το (2,1) στο (7,4).} \\
            }
        }
    }
     \\ \\

    \theutterance \stepcounter{utterance}  
    & & & \multicolumn{2}{p{0.3\linewidth}}{
        \cellcolor[rgb]{0.9,0.9,0.9}{
            \makecell[{{p{\linewidth}}}]{
                \texttt{\tiny{[GM$|$GM]}}
                \texttt{ΠΕΣ: Έλεγξα τις προτάσεις σου στον πίνακά μου: (3,7) είναι άδειο, (1,4) είναι άδειο, (7,4) είναι άδειο. Συμφωνώ με C στο (3,7), L στο (1,4), P στο (7,4). Να ξεκινήσω τις μετακινήσεις; Αν συμφωνείς, πες μου να ξεκινήσουμε και εγώ θα μετακινήσω πρώτα το P από το (2,1) στο (7,4).} \\
            }
        }
    }
    & & \\ \\

    \theutterance \stepcounter{utterance}  
    & & \multicolumn{4}{p{0.6\linewidth}}{
        \cellcolor[rgb]{0.9,0.9,0.9}{
            \makecell[{{p{\linewidth}}}]{
                \texttt{\tiny{[P1$\langle$GM]}}
                \texttt{Το μήνυμά σας μεταφέρθηκε στον άλλο παίκτη.} \\
\\ 
\\ 
\texttt{Αυτή τη στιγμή παίζετε τον γύρο 3 από συνολικά 12.} \\
\texttt{Έχετε συσσωρεύσει συλλογικά 1 από 8 ποινές.} \\
\texttt{Ο άλλος παίκτης έστειλε αυτό το μήνυμα:} \\
\\ 
\\ 
\texttt{"Έλεγξα τις προτάσεις σου στον πίνακά μου: (3,7) είναι άδειο, (1,4) είναι άδειο, (7,4) είναι άδειο. Συμφωνώ με C στο (3,7), L στο (1,4), P στο (7,4). Να ξεκινήσω τις μετακινήσεις; Αν συμφωνείς, πες μου να ξεκινήσουμε και εγώ θα μετακινήσω πρώτα το P από το (2,1) στο (7,4)."} \\
\\ 
\\ 
\texttt{Ποια είναι η επόμενη εντολή σας;} \\
            }
        }
    }
    & \\ \\

    \theutterance \stepcounter{utterance}  
    & \multicolumn{4}{p{0.6\linewidth}}{
        \cellcolor[rgb]{0.9,0.9,0.9}{
            \makecell[{{p{\linewidth}}}]{
                \texttt{\tiny{[P1$\rangle$GM]}}
                \texttt{ΠΕΣ: Συμφωνώ! Ας ξεκινήσουμε. Θα μετακινήσω πρώτα το P από το (6,3) στο (7,4). Πάμε!} \\
            }
        }
    }
    & & \\ \\

    \theutterance \stepcounter{utterance}  
    & & & \multicolumn{2}{p{0.3\linewidth}}{
        \cellcolor[rgb]{0.9,0.9,0.9}{
            \makecell[{{p{\linewidth}}}]{
                \texttt{\tiny{[GM$|$GM]}}
                \texttt{ΠΕΣ: Συμφωνώ! Ας ξεκινήσουμε. Θα μετακινήσω πρώτα το P από το (6,3) στο (7,4). Πάμε!} \\
            }
        }
    }
    & & \\ \\

    \theutterance \stepcounter{utterance}  
    & & \multicolumn{4}{p{0.6\linewidth}}{
        \cellcolor[rgb]{0.9,0.9,0.9}{
            \makecell[{{p{\linewidth}}}]{
                \texttt{\tiny{[GM$\rangle$P2]}}
                \texttt{Το μήνυμά σας μεταφέρθηκε στον άλλο παίκτη.} \\
\\ 
\\ 
\texttt{Αυτή τη στιγμή παίζετε τον γύρο 4 από συνολικά 12.} \\
\texttt{Έχετε συσσωρεύσει συλλογικά 1 από 8 ποινές.} \\
\texttt{Ο άλλος παίκτης έστειλε αυτό το μήνυμα:} \\
\\ 
\\ 
\texttt{"Συμφωνώ! Ας ξεκινήσουμε. Θα μετακινήσω πρώτα το P από το (6,3) στο (7,4). Πάμε!"} \\
\\ 
\\ 
\texttt{Ποια είναι η επόμενη εντολή σας;} \\
            }
        }
    }
    & \\ \\

    \theutterance \stepcounter{utterance}  
    & & & \multicolumn{4}{p{0.6\linewidth}}{
        \cellcolor[rgb]{0.9,0.9,0.9}{
            \makecell[{{p{\linewidth}}}]{
                \texttt{\tiny{[GM$\langle$P2]}}
                \texttt{ΚΙΝΗΣΗ: P, (7,4)} \\
            }
        }
    }
     \\ \\

    \theutterance \stepcounter{utterance}  
    & & & \multicolumn{2}{p{0.3\linewidth}}{
        \cellcolor[rgb]{0.9,0.9,0.9}{
            \makecell[{{p{\linewidth}}}]{
                \texttt{\tiny{[GM$|$GM]}}
                \texttt{ΚΙΝΗΣΗ: P, (7,4)} \\
            }
        }
    }
    & & \\ \\

    \theutterance \stepcounter{utterance}  
    & & & \multicolumn{2}{p{0.3\linewidth}}{
        \cellcolor[rgb]{0.9,0.9,0.9}{
            \makecell[{{p{\linewidth}}}]{
                \texttt{\tiny{[GM$|$GM]}}
                \texttt{Ποινή: Το (7,4) δεν είναι κενό, αλλά περιέχει '─'.} \\
\\ 
\texttt{Έχετε συσσωρεύσει συλλογικά 2 από 8 ποινές.} \\
\texttt{Παρακαλώ δοκιμάστε ξανά!} \\
            }
        }
    }
    & & \\ \\

    \theutterance \stepcounter{utterance}  
    & & \multicolumn{4}{p{0.6\linewidth}}{
        \cellcolor[rgb]{0.9,0.9,0.9}{
            \makecell[{{p{\linewidth}}}]{
                \texttt{\tiny{[GM$\rangle$P2]}}
                \texttt{Ποινή: Το (7,4) δεν είναι κενό, αλλά περιέχει '─'.} \\
\\ 
\texttt{Έχετε συσσωρεύσει συλλογικά 2 από 8 ποινές.} \\
\texttt{Παρακαλώ δοκιμάστε ξανά!} \\
            }
        }
    }
    & \\ \\

    \theutterance \stepcounter{utterance}  
    & & & \multicolumn{4}{p{0.6\linewidth}}{
        \cellcolor[rgb]{0.9,0.9,0.9}{
            \makecell[{{p{\linewidth}}}]{
                \texttt{\tiny{[GM$\langle$P2]}}
                \texttt{ΠΕΣ: Στον πίνακά μου το (7,4) περιέχει '─', δεν είναι άδειο! Χρειάζεται να ξαναελέγξουμε. Μήπως έκανα λάθος πριν; Ας ξαναδούμε τα κοινά άδεια κελιά. Στον δικό μου πίνακα, το (7,4) έχει σύμβολο. Μπορείς να επιβεβαιώσεις τα δικά σου άδεια στο (7,4); Ας βρούμε εναλλακτική θέση για το P. Τι θα έλεγες για P στο (7,5); Είναι άδειο στον δικό μου πίνακα.} \\
            }
        }
    }
     \\ \\

    \theutterance \stepcounter{utterance}  
    & & & \multicolumn{2}{p{0.3\linewidth}}{
        \cellcolor[rgb]{0.9,0.9,0.9}{
            \makecell[{{p{\linewidth}}}]{
                \texttt{\tiny{[GM$|$GM]}}
                \texttt{ΠΕΣ: Στον πίνακά μου το (7,4) περιέχει '─', δεν είναι άδειο! Χρειάζεται να ξαναελέγξουμε. Μήπως έκανα λάθος πριν; Ας ξαναδούμε τα κοινά άδεια κελιά. Στον δικό μου πίνακα, το (7,4) έχει σύμβολο. Μπορείς να επιβεβαιώσεις τα δικά σου άδεια στο (7,4); Ας βρούμε εναλλακτική θέση για το P. Τι θα έλεγες για P στο (7,5); Είναι άδειο στον δικό μου πίνακα.} \\
            }
        }
    }
    & & \\ \\

    \theutterance \stepcounter{utterance}  
    & & \multicolumn{4}{p{0.6\linewidth}}{
        \cellcolor[rgb]{0.9,0.9,0.9}{
            \makecell[{{p{\linewidth}}}]{
                \texttt{\tiny{[P1$\langle$GM]}}
                \texttt{Το μήνυμά σας μεταφέρθηκε στον άλλο παίκτη.} \\
\\ 
\\ 
\texttt{Αυτή τη στιγμή παίζετε τον γύρο 4 από συνολικά 12.} \\
\texttt{Έχετε συσσωρεύσει συλλογικά 2 από 8 ποινές.} \\
\texttt{Ο άλλος παίκτης έστειλε αυτό το μήνυμα:} \\
\\ 
\\ 
\texttt{"Στον πίνακά μου το (7,4) περιέχει '─', δεν είναι άδειο! Χρειάζεται να ξαναελέγξουμε. Μήπως έκανα λάθος πριν; Ας ξαναδούμε τα κοινά άδεια κελιά. Στον δικό μου πίνακα, το (7,4) έχει σύμβολο. Μπορείς να επιβεβαιώσεις τα δικά σου άδεια στο (7,4); Ας βρούμε εναλλακτική θέση για το P. Τι θα έλεγες για P στο (7,5); Είναι άδειο στον δικό μου πίνακα."} \\
\\ 
\\ 
\texttt{Ποια είναι η επόμενη εντολή σας;} \\
            }
        }
    }
    & \\ \\

    \theutterance \stepcounter{utterance}  
    & \multicolumn{4}{p{0.6\linewidth}}{
        \cellcolor[rgb]{0.9,0.9,0.9}{
            \makecell[{{p{\linewidth}}}]{
                \texttt{\tiny{[P1$\rangle$GM]}}
                \texttt{ΠΕΣ: Έχεις δίκιο, λυπάμαι! Στον δικό μου πίνακα το (7,4) είναι άδειο, αλλά πρέπει να είναι κοινό. Το (7,5) είναι επίσης άδειο στον δικό μου. Αλλά πρόσεξε: στον δικό μου πίνακα, το (3,7) που προτείναμε για C είναι άδειο, το (1,4) για L είναι άδειο. Ας επιβεβαιώσουμε αυτά πρώτα. Για P, ας δοκιμάσουμε το (7,5). Είναι άδειο στον δικό μου. Το επιβεβαιώνεις; Νέοι στόχοι: C(3,7), L(1,4), P(7,5).} \\
            }
        }
    }
    & & \\ \\

    \theutterance \stepcounter{utterance}  
    & & & \multicolumn{2}{p{0.3\linewidth}}{
        \cellcolor[rgb]{0.9,0.9,0.9}{
            \makecell[{{p{\linewidth}}}]{
                \texttt{\tiny{[GM$|$GM]}}
                \texttt{ΠΕΣ: Έχεις δίκιο, λυπάμαι! Στον δικό μου πίνακα το (7,4) είναι άδειο, αλλά πρέπει να είναι κοινό. Το (7,5) είναι επίσης άδειο στον δικό μου. Αλλά πρόσεξε: στον δικό μου πίνακα, το (3,7) που προτείναμε για C είναι άδειο, το (1,4) για L είναι άδειο. Ας επιβεβαιώσουμε αυτά πρώτα. Για P, ας δοκιμάσουμε το (7,5). Είναι άδειο στον δικό μου. Το επιβεβαιώνεις; Νέοι στόχοι: C(3,7), L(1,4), P(7,5).} \\
            }
        }
    }
    & & \\ \\

    \theutterance \stepcounter{utterance}  
    & & \multicolumn{4}{p{0.6\linewidth}}{
        \cellcolor[rgb]{0.9,0.9,0.9}{
            \makecell[{{p{\linewidth}}}]{
                \texttt{\tiny{[GM$\rangle$P2]}}
                \texttt{Το μήνυμά σας μεταφέρθηκε στον άλλο παίκτη.} \\
\\ 
\\ 
\texttt{Αυτή τη στιγμή παίζετε τον γύρο 5 από συνολικά 12.} \\
\texttt{Έχετε συσσωρεύσει συλλογικά 2 από 8 ποινές.} \\
\texttt{Ο άλλος παίκτης έστειλε αυτό το μήνυμα:} \\
\\ 
\\ 
\texttt{"Έχεις δίκιο, λυπάμαι! Στον δικό μου πίνακα το (7,4) είναι άδειο, αλλά πρέπει να είναι κοινό. Το (7,5) είναι επίσης άδειο στον δικό μου. Αλλά πρόσεξε: στον δικό μου πίνακα, το (3,7) που προτείναμε για C είναι άδειο, το (1,4) για L είναι άδειο. Ας επιβεβαιώσουμε αυτά πρώτα. Για P, ας δοκιμάσουμε το (7,5). Είναι άδειο στον δικό μου. Το επιβεβαιώνεις; Νέοι στόχοι: C(3,7), L(1,4), P(7,5)."} \\
\\ 
\\ 
\texttt{Ποια είναι η επόμενη εντολή σας;} \\
            }
        }
    }
    & \\ \\

    \theutterance \stepcounter{utterance}  
    & & & \multicolumn{4}{p{0.6\linewidth}}{
        \cellcolor[rgb]{0.9,0.9,0.9}{
            \makecell[{{p{\linewidth}}}]{
                \texttt{\tiny{[GM$\langle$P2]}}
                \texttt{ΠΕΣ: Επιβεβαιώνω: (3,7) είναι άδειο, (1,4) είναι άδειο, (7,5) είναι άδειο στον πίνακά μου. Συμφωνώ: C(3,7), L(1,4), P(7,5). Ξεκινάω μετακινήσεις! Πρώτα μετακινώ το P από (2,1) στο (7,5).} \\
            }
        }
    }
     \\ \\

    \theutterance \stepcounter{utterance}  
    & & & \multicolumn{2}{p{0.3\linewidth}}{
        \cellcolor[rgb]{0.9,0.9,0.9}{
            \makecell[{{p{\linewidth}}}]{
                \texttt{\tiny{[GM$|$GM]}}
                \texttt{ΠΕΣ: Επιβεβαιώνω: (3,7) είναι άδειο, (1,4) είναι άδειο, (7,5) είναι άδειο στον πίνακά μου. Συμφωνώ: C(3,7), L(1,4), P(7,5). Ξεκινάω μετακινήσεις! Πρώτα μετακινώ το P από (2,1) στο (7,5).} \\
            }
        }
    }
    & & \\ \\

    \theutterance \stepcounter{utterance}  
    & & \multicolumn{4}{p{0.6\linewidth}}{
        \cellcolor[rgb]{0.9,0.9,0.9}{
            \makecell[{{p{\linewidth}}}]{
                \texttt{\tiny{[P1$\langle$GM]}}
                \texttt{Το μήνυμά σας μεταφέρθηκε στον άλλο παίκτη.} \\
\\ 
\\ 
\texttt{Αυτή τη στιγμή παίζετε τον γύρο 5 από συνολικά 12.} \\
\texttt{Έχετε συσσωρεύσει συλλογικά 2 από 8 ποινές.} \\
\texttt{Ο άλλος παίκτης έστειλε αυτό το μήνυμα:} \\
\\ 
\\ 
\texttt{"Επιβεβαιώνω: (3,7) είναι άδειο, (1,4) είναι άδειο, (7,5) είναι άδειο στον πίνακά μου. Συμφωνώ: C(3,7), L(1,4), P(7,5). Ξεκινάω μετακινήσεις! Πρώτα μετακινώ το P από (2,1) στο (7,5)."} \\
\\ 
\\ 
\texttt{Ποια είναι η επόμενη εντολή σας;} \\
            }
        }
    }
    & \\ \\

    \theutterance \stepcounter{utterance}  
    & \multicolumn{4}{p{0.6\linewidth}}{
        \cellcolor[rgb]{0.9,0.9,0.9}{
            \makecell[{{p{\linewidth}}}]{
                \texttt{\tiny{[P1$\rangle$GM]}}
                \texttt{ΚΙΝΗΣΗ: P, (7,5)} \\
            }
        }
    }
    & & \\ \\

    \theutterance \stepcounter{utterance}  
    & & & \multicolumn{2}{p{0.3\linewidth}}{
        \cellcolor[rgb]{0.9,0.9,0.9}{
            \makecell[{{p{\linewidth}}}]{
                \texttt{\tiny{[GM$|$GM]}}
                \texttt{ΚΙΝΗΣΗ: P, (7,5)} \\
            }
        }
    }
    & & \\ \\

    \theutterance \stepcounter{utterance}  
    & & & \multicolumn{2}{p{0.3\linewidth}}{
        \cellcolor[rgb]{0.9,0.9,0.9}{
            \makecell[{{p{\linewidth}}}]{
                \texttt{\tiny{[GM$|$GM]}}
                \texttt{P moved to (7, 5)} \\
            }
        }
    }
    & & \\ \\

    \theutterance \stepcounter{utterance}  
    & & \multicolumn{4}{p{0.6\linewidth}}{
        \cellcolor[rgb]{0.9,0.9,0.9}{
            \makecell[{{p{\linewidth}}}]{
                \texttt{\tiny{[GM$\rangle$P2]}}
                \texttt{Το μήνυμά σας μεταφέρθηκε στον άλλο παίκτη.} \\
\\ 
\\ 
\texttt{Αυτή τη στιγμή παίζετε τον γύρο 6 από συνολικά 12.} \\
\texttt{Έχετε συσσωρεύσει συλλογικά 2 από 8 ποινές.} \\
\texttt{Ο άλλος παίκτης μετακίνησε ένα αντικείμενο στο πλέγμα του.} \\
\\ 
\\ 
\texttt{Ποια είναι η επόμενη εντολή σας;} \\
            }
        }
    }
    & \\ \\

    \theutterance \stepcounter{utterance}  
    & & & \multicolumn{4}{p{0.6\linewidth}}{
        \cellcolor[rgb]{0.9,0.9,0.9}{
            \makecell[{{p{\linewidth}}}]{
                \texttt{\tiny{[GM$\langle$P2]}}
                \texttt{ΚΙΝΗΣΗ: P, (7,5)} \\
            }
        }
    }
     \\ \\

    \theutterance \stepcounter{utterance}  
    & & & \multicolumn{2}{p{0.3\linewidth}}{
        \cellcolor[rgb]{0.9,0.9,0.9}{
            \makecell[{{p{\linewidth}}}]{
                \texttt{\tiny{[GM$|$GM]}}
                \texttt{ΚΙΝΗΣΗ: P, (7,5)} \\
            }
        }
    }
    & & \\ \\

    \theutterance \stepcounter{utterance}  
    & & & \multicolumn{2}{p{0.3\linewidth}}{
        \cellcolor[rgb]{0.9,0.9,0.9}{
            \makecell[{{p{\linewidth}}}]{
                \texttt{\tiny{[GM$|$GM]}}
                \texttt{P moved to (7, 5)} \\
            }
        }
    }
    & & \\ \\

    \theutterance \stepcounter{utterance}  
    & & \multicolumn{4}{p{0.6\linewidth}}{
        \cellcolor[rgb]{0.9,0.9,0.9}{
            \makecell[{{p{\linewidth}}}]{
                \texttt{\tiny{[P1$\langle$GM]}}
                \texttt{Μεταφέρθηκε επιτυχώς το 'P' στο (7,5). Το ενημερωμένο πλέγμα σας μοιάζει ως εξής:} \\
\texttt{```} \\
\texttt{ 1234567} \\
\texttt{╔══╤═╤══╗} \\
\texttt{║◌◌│◌│◌◌║ 1} \\
\texttt{║◌◌└─┼──╢ 2} \\
\texttt{║◌◌◌◌│◌◌║ 3} \\
\texttt{║◌┌─┬┴┬─╢ 4} \\
\texttt{╟─┤◌│◌│P║ 5} \\
\texttt{║◌└┬┴┬┘L║ 6} \\
\texttt{║◌◌│C│◌◌║ 7} \\
\texttt{╚══╧═╧══╝} \\
\texttt{```} \\
\texttt{Αυτή τη στιγμή παίζετε τον γύρο 6 από συνολικά 12.} \\
\texttt{Έχετε συσσωρεύσει συλλογικά 2 από 8 ποινές.} \\
\texttt{Ο άλλος παίκτης μετακίνησε ένα αντικείμενο στο πλέγμα του.} \\
\\ 
\\ 
\texttt{Ποια είναι η επόμενη εντολή σας;} \\
            }
        }
    }
    & \\ \\

    \theutterance \stepcounter{utterance}  
    & \multicolumn{4}{p{0.6\linewidth}}{
        \cellcolor[rgb]{0.9,0.9,0.9}{
            \makecell[{{p{\linewidth}}}]{
                \texttt{\tiny{[P1$\rangle$GM]}}
                \texttt{ΠΕΣ: Ωραία! Το P μου μετακινήθηκε στο (7,5). Τώρα ας μετακινήσουμε τα άλλα. Θα μετακινήσω το C από το (4,7) στο (3,7). Εσύ;} \\
            }
        }
    }
    & & \\ \\

    \theutterance \stepcounter{utterance}  
    & & & \multicolumn{2}{p{0.3\linewidth}}{
        \cellcolor[rgb]{0.9,0.9,0.9}{
            \makecell[{{p{\linewidth}}}]{
                \texttt{\tiny{[GM$|$GM]}}
                \texttt{ΠΕΣ: Ωραία! Το P μου μετακινήθηκε στο (7,5). Τώρα ας μετακινήσουμε τα άλλα. Θα μετακινήσω το C από το (4,7) στο (3,7). Εσύ;} \\
            }
        }
    }
    & & \\ \\

    \theutterance \stepcounter{utterance}  
    & & \multicolumn{4}{p{0.6\linewidth}}{
        \cellcolor[rgb]{0.9,0.9,0.9}{
            \makecell[{{p{\linewidth}}}]{
                \texttt{\tiny{[GM$\rangle$P2]}}
                \texttt{Μεταφέρθηκε επιτυχώς το 'P' στο (7,5). Το ενημερωμένο πλέγμα σας μοιάζει ως εξής:} \\
\texttt{```} \\
\texttt{ 1234567} \\
\texttt{╔══╤═╤══╗} \\
\texttt{║◌◌│◌│◌◌║ 1} \\
\texttt{║◌L└─┼──╢ 2} \\
\texttt{║◌◌◌◌│◌◌║ 3} \\
\texttt{║◌┌─┬┴┬─╢ 4} \\
\texttt{╟─┤◌│◌│P║ 5} \\
\texttt{║◌└┬┴┬┘◌║ 6} \\
\texttt{║◌◌│◌│◌C║ 7} \\
\texttt{╚══╧═╧══╝} \\
\texttt{```} \\
\texttt{Αυτή τη στιγμή παίζετε τον γύρο 7 από συνολικά 12.} \\
\texttt{Έχετε συσσωρεύσει συλλογικά 2 από 8 ποινές.} \\
\texttt{Ο άλλος παίκτης έστειλε αυτό το μήνυμα:} \\
\\ 
\\ 
\texttt{"Ωραία! Το P μου μετακινήθηκε στο (7,5). Τώρα ας μετακινήσουμε τα άλλα. Θα μετακινήσω το C από το (4,7) στο (3,7). Εσύ;"} \\
\\ 
\\ 
\texttt{Ποια είναι η επόμενη εντολή σας;} \\
            }
        }
    }
    & \\ \\

    \theutterance \stepcounter{utterance}  
    & & & \multicolumn{4}{p{0.6\linewidth}}{
        \cellcolor[rgb]{0.9,0.9,0.9}{
            \makecell[{{p{\linewidth}}}]{
                \texttt{\tiny{[GM$\langle$P2]}}
                \texttt{ΚΙΝΗΣΗ: C, (3,7)} \\
            }
        }
    }
     \\ \\

    \theutterance \stepcounter{utterance}  
    & & & \multicolumn{2}{p{0.3\linewidth}}{
        \cellcolor[rgb]{0.9,0.9,0.9}{
            \makecell[{{p{\linewidth}}}]{
                \texttt{\tiny{[GM$|$GM]}}
                \texttt{ΚΙΝΗΣΗ: C, (3,7)} \\
            }
        }
    }
    & & \\ \\

    \theutterance \stepcounter{utterance}  
    & & & \multicolumn{2}{p{0.3\linewidth}}{
        \cellcolor[rgb]{0.9,0.9,0.9}{
            \makecell[{{p{\linewidth}}}]{
                \texttt{\tiny{[GM$|$GM]}}
                \texttt{Ποινή: Το (3,7) δεν είναι κενό, αλλά περιέχει '│'.} \\
\\ 
\texttt{Έχετε συσσωρεύσει συλλογικά 3 από 8 ποινές.} \\
\texttt{Παρακαλώ δοκιμάστε ξανά!} \\
            }
        }
    }
    & & \\ \\

    \theutterance \stepcounter{utterance}  
    & & \multicolumn{4}{p{0.6\linewidth}}{
        \cellcolor[rgb]{0.9,0.9,0.9}{
            \makecell[{{p{\linewidth}}}]{
                \texttt{\tiny{[GM$\rangle$P2]}}
                \texttt{Ποινή: Το (3,7) δεν είναι κενό, αλλά περιέχει '│'.} \\
\\ 
\texttt{Έχετε συσσωρεύσει συλλογικά 3 από 8 ποινές.} \\
\texttt{Παρακαλώ δοκιμάστε ξανά!} \\
            }
        }
    }
    & \\ \\

    \theutterance \stepcounter{utterance}  
    & & & \multicolumn{4}{p{0.6\linewidth}}{
        \cellcolor[rgb]{0.9,0.9,0.9}{
            \makecell[{{p{\linewidth}}}]{
                \texttt{\tiny{[GM$\langle$P2]}}
                \texttt{ΠΕΣ: Στον πίνακά μου το (3,7) έχει '│', δεν είναι άδειο! Πρέπει να το ξαναελέγξουμε. Στη σειρά 7, τα κελιά μου είναι: (1,7) άδειο, (2,7) άδειο, (3,7) '│', (4,7) άδειο, (5,7) άδειο, (6,7) άδειο, (7,7) C. Μήπως μπερδεύτηκα πριν; Ας προτείνω εναλλακτική: το (4,7) είναι άδειο στον δικό μου. Είναι άδειο και στον δικό σου; Αν ναι, προτείνω νέο στόχο για C: (4,7). Επίσης ας επιβεβαιώσουμε το L στο (1,4) {-} είναι όντως άδειο στον δικό μου.} \\
            }
        }
    }
     \\ \\

    \theutterance \stepcounter{utterance}  
    & & & \multicolumn{2}{p{0.3\linewidth}}{
        \cellcolor[rgb]{0.9,0.9,0.9}{
            \makecell[{{p{\linewidth}}}]{
                \texttt{\tiny{[GM$|$GM]}}
                \texttt{ΠΕΣ: Στον πίνακά μου το (3,7) έχει '│', δεν είναι άδειο! Πρέπει να το ξαναελέγξουμε. Στη σειρά 7, τα κελιά μου είναι: (1,7) άδειο, (2,7) άδειο, (3,7) '│', (4,7) άδειο, (5,7) άδειο, (6,7) άδειο, (7,7) C. Μήπως μπερδεύτηκα πριν; Ας προτείνω εναλλακτική: το (4,7) είναι άδειο στον δικό μου. Είναι άδειο και στον δικό σου; Αν ναι, προτείνω νέο στόχο για C: (4,7). Επίσης ας επιβεβαιώσουμε το L στο (1,4) {-} είναι όντως άδειο στον δικό μου.} \\
            }
        }
    }
    & & \\ \\

    \theutterance \stepcounter{utterance}  
    & & \multicolumn{4}{p{0.6\linewidth}}{
        \cellcolor[rgb]{0.9,0.9,0.9}{
            \makecell[{{p{\linewidth}}}]{
                \texttt{\tiny{[P1$\langle$GM]}}
                \texttt{Το μήνυμά σας μεταφέρθηκε στον άλλο παίκτη.} \\
\\ 
\\ 
\texttt{Αυτή τη στιγμή παίζετε τον γύρο 7 από συνολικά 12.} \\
\texttt{Έχετε συσσωρεύσει συλλογικά 3 από 8 ποινές.} \\
\texttt{Ο άλλος παίκτης έστειλε αυτό το μήνυμα:} \\
\\ 
\\ 
\texttt{"Στον πίνακά μου το (3,7) έχει '│', δεν είναι άδειο! Πρέπει να το ξαναελέγξουμε. Στη σειρά 7, τα κελιά μου είναι: (1,7) άδειο, (2,7) άδειο, (3,7) '│', (4,7) άδειο, (5,7) άδειο, (6,7) άδειο, (7,7) C. Μήπως μπερδεύτηκα πριν; Ας προτείνω εναλλακτική: το (4,7) είναι άδειο στον δικό μου. Είναι άδειο και στον δικό σου; Αν ναι, προτείνω νέο στόχο για C: (4,7). Επίσης ας επιβεβαιώσουμε το L στο (1,4) {-} είναι όντως άδειο στον δικό μου."} \\
\\ 
\\ 
\texttt{Ποια είναι η επόμενη εντολή σας;} \\
            }
        }
    }
    & \\ \\

    \theutterance \stepcounter{utterance}  
    & \multicolumn{4}{p{0.6\linewidth}}{
        \cellcolor[rgb]{0.9,0.9,0.9}{
            \makecell[{{p{\linewidth}}}]{
                \texttt{\tiny{[P1$\rangle$GM]}}
                \texttt{ΠΕΣ: Περίμενε, στον δικό μου πίνακα το (4,7) έχει ήδη το C! Δεν είναι άδειο, είναι η αρχική θέση του C. Άρα δεν γίνεται να είναι στόχος. Ας βρούμε άλλη θέση. Το (5,7) είναι άδειο στον δικό μου, το ίδιο και το (6,7). Τι θα έλεγες για C στο (5,7); Είναι άδειο στον δικό σου; Επίσης επιβεβαιώνω ότι το (1,4) είναι άδειο στον δικό μου για το L.} \\
            }
        }
    }
    & & \\ \\

    \theutterance \stepcounter{utterance}  
    & & & \multicolumn{2}{p{0.3\linewidth}}{
        \cellcolor[rgb]{0.9,0.9,0.9}{
            \makecell[{{p{\linewidth}}}]{
                \texttt{\tiny{[GM$|$GM]}}
                \texttt{ΠΕΣ: Περίμενε, στον δικό μου πίνακα το (4,7) έχει ήδη το C! Δεν είναι άδειο, είναι η αρχική θέση του C. Άρα δεν γίνεται να είναι στόχος. Ας βρούμε άλλη θέση. Το (5,7) είναι άδειο στον δικό μου, το ίδιο και το (6,7). Τι θα έλεγες για C στο (5,7); Είναι άδειο στον δικό σου; Επίσης επιβεβαιώνω ότι το (1,4) είναι άδειο στον δικό μου για το L.} \\
            }
        }
    }
    & & \\ \\

    \theutterance \stepcounter{utterance}  
    & & \multicolumn{4}{p{0.6\linewidth}}{
        \cellcolor[rgb]{0.9,0.9,0.9}{
            \makecell[{{p{\linewidth}}}]{
                \texttt{\tiny{[GM$\rangle$P2]}}
                \texttt{Το μήνυμά σας μεταφέρθηκε στον άλλο παίκτη.} \\
\\ 
\\ 
\texttt{Αυτή τη στιγμή παίζετε τον γύρο 8 από συνολικά 12.} \\
\texttt{Έχετε συσσωρεύσει συλλογικά 3 από 8 ποινές.} \\
\texttt{Ο άλλος παίκτης έστειλε αυτό το μήνυμα:} \\
\\ 
\\ 
\texttt{"Περίμενε, στον δικό μου πίνακα το (4,7) έχει ήδη το C! Δεν είναι άδειο, είναι η αρχική θέση του C. Άρα δεν γίνεται να είναι στόχος. Ας βρούμε άλλη θέση. Το (5,7) είναι άδειο στον δικό μου, το ίδιο και το (6,7). Τι θα έλεγες για C στο (5,7); Είναι άδειο στον δικό σου; Επίσης επιβεβαιώνω ότι το (1,4) είναι άδειο στον δικό μου για το L."} \\
\\ 
\\ 
\texttt{Ποια είναι η επόμενη εντολή σας;} \\
            }
        }
    }
    & \\ \\

    \theutterance \stepcounter{utterance}  
    & & & \multicolumn{4}{p{0.6\linewidth}}{
        \cellcolor[rgb]{0.9,0.9,0.9}{
            \makecell[{{p{\linewidth}}}]{
                \texttt{\tiny{[GM$\langle$P2]}}
                \texttt{ΠΕΣ: Στον πίνακά μου το (5,7) είναι άδειο και το (6,7) είναι άδειο. Συμφωνώ για C στο (5,7). Οπότε νέοι στόχοι: C(5,7), L(1,4), P(7,5). Το (1,4) είναι άδειο, το (7,5) έχει ήδη το P. Πάω να μετακινήσω το C από (7,7) στο (5,7).} \\
            }
        }
    }
     \\ \\

    \theutterance \stepcounter{utterance}  
    & & & \multicolumn{2}{p{0.3\linewidth}}{
        \cellcolor[rgb]{0.9,0.9,0.9}{
            \makecell[{{p{\linewidth}}}]{
                \texttt{\tiny{[GM$|$GM]}}
                \texttt{ΠΕΣ: Στον πίνακά μου το (5,7) είναι άδειο και το (6,7) είναι άδειο. Συμφωνώ για C στο (5,7). Οπότε νέοι στόχοι: C(5,7), L(1,4), P(7,5). Το (1,4) είναι άδειο, το (7,5) έχει ήδη το P. Πάω να μετακινήσω το C από (7,7) στο (5,7).} \\
            }
        }
    }
    & & \\ \\

    \theutterance \stepcounter{utterance}  
    & & \multicolumn{4}{p{0.6\linewidth}}{
        \cellcolor[rgb]{0.9,0.9,0.9}{
            \makecell[{{p{\linewidth}}}]{
                \texttt{\tiny{[P1$\langle$GM]}}
                \texttt{Το μήνυμά σας μεταφέρθηκε στον άλλο παίκτη.} \\
\\ 
\\ 
\texttt{Αυτή τη στιγμή παίζετε τον γύρο 8 από συνολικά 12.} \\
\texttt{Έχετε συσσωρεύσει συλλογικά 3 από 8 ποινές.} \\
\texttt{Ο άλλος παίκτης έστειλε αυτό το μήνυμα:} \\
\\ 
\\ 
\texttt{"Στον πίνακά μου το (5,7) είναι άδειο και το (6,7) είναι άδειο. Συμφωνώ για C στο (5,7). Οπότε νέοι στόχοι: C(5,7), L(1,4), P(7,5). Το (1,4) είναι άδειο, το (7,5) έχει ήδη το P. Πάω να μετακινήσω το C από (7,7) στο (5,7)."} \\
\\ 
\\ 
\texttt{Ποια είναι η επόμενη εντολή σας;} \\
            }
        }
    }
    & \\ \\

    \theutterance \stepcounter{utterance}  
    & \multicolumn{4}{p{0.6\linewidth}}{
        \cellcolor[rgb]{0.9,0.9,0.9}{
            \makecell[{{p{\linewidth}}}]{
                \texttt{\tiny{[P1$\rangle$GM]}}
                \texttt{ΚΙΝΗΣΗ: C, (5,7)} \\
            }
        }
    }
    & & \\ \\

    \theutterance \stepcounter{utterance}  
    & & & \multicolumn{2}{p{0.3\linewidth}}{
        \cellcolor[rgb]{0.9,0.9,0.9}{
            \makecell[{{p{\linewidth}}}]{
                \texttt{\tiny{[GM$|$GM]}}
                \texttt{ΚΙΝΗΣΗ: C, (5,7)} \\
            }
        }
    }
    & & \\ \\

    \theutterance \stepcounter{utterance}  
    & & & \multicolumn{2}{p{0.3\linewidth}}{
        \cellcolor[rgb]{0.9,0.9,0.9}{
            \makecell[{{p{\linewidth}}}]{
                \texttt{\tiny{[GM$|$GM]}}
                \texttt{Ποινή: Το (5,7) δεν είναι κενό, αλλά περιέχει '│'.} \\
\\ 
\texttt{Έχετε συσσωρεύσει συλλογικά 4 από 8 ποινές.} \\
\texttt{Παρακαλώ δοκιμάστε ξανά!} \\
            }
        }
    }
    & & \\ \\

    \theutterance \stepcounter{utterance}  
    & & \multicolumn{4}{p{0.6\linewidth}}{
        \cellcolor[rgb]{0.9,0.9,0.9}{
            \makecell[{{p{\linewidth}}}]{
                \texttt{\tiny{[P1$\langle$GM]}}
                \texttt{Ποινή: Το (5,7) δεν είναι κενό, αλλά περιέχει '│'.} \\
\\ 
\texttt{Έχετε συσσωρεύσει συλλογικά 4 από 8 ποινές.} \\
\texttt{Παρακαλώ δοκιμάστε ξανά!} \\
            }
        }
    }
    & \\ \\

    \theutterance \stepcounter{utterance}  
    & \multicolumn{4}{p{0.6\linewidth}}{
        \cellcolor[rgb]{0.9,0.9,0.9}{
            \makecell[{{p{\linewidth}}}]{
                \texttt{\tiny{[P1$\rangle$GM]}}
                \texttt{ΠΕΣ: Έλαβα ποινή! Στον πίνακά μου το (5,7) έχει '│', δεν είναι άδειο. Λάθος μου που δεν το έλεγξα. Στη σειρά 7 του πίνακά μου: (1,7) άδειο, (2,7) άδειο, (3,7) άδειο, (4,7) C, (5,7) '│', (6,7) άδειο, (7,7) άδειο. Άρα για C, τα άδεια είναι (1,7), (2,7), (3,7), (6,7), (7,7). Εσύ είπες ότι (1,7), (2,7), (6,7) είναι άδεια. Ας διαλέξουμε ένα από αυτά. Προτείνω C στο (6,7). Είναι άδειο στον δικό σου;} \\
            }
        }
    }
    & & \\ \\

    \theutterance \stepcounter{utterance}  
    & & & \multicolumn{2}{p{0.3\linewidth}}{
        \cellcolor[rgb]{0.9,0.9,0.9}{
            \makecell[{{p{\linewidth}}}]{
                \texttt{\tiny{[GM$|$GM]}}
                \texttt{ΠΕΣ: Έλαβα ποινή! Στον πίνακά μου το (5,7) έχει '│', δεν είναι άδειο. Λάθος μου που δεν το έλεγξα. Στη σειρά 7 του πίνακά μου: (1,7) άδειο, (2,7) άδειο, (3,7) άδειο, (4,7) C, (5,7) '│', (6,7) άδειο, (7,7) άδειο. Άρα για C, τα άδεια είναι (1,7), (2,7), (3,7), (6,7), (7,7). Εσύ είπες ότι (1,7), (2,7), (6,7) είναι άδεια. Ας διαλέξουμε ένα από αυτά. Προτείνω C στο (6,7). Είναι άδειο στον δικό σου;} \\
            }
        }
    }
    & & \\ \\

    \theutterance \stepcounter{utterance}  
    & & \multicolumn{4}{p{0.6\linewidth}}{
        \cellcolor[rgb]{0.9,0.9,0.9}{
            \makecell[{{p{\linewidth}}}]{
                \texttt{\tiny{[GM$\rangle$P2]}}
                \texttt{Το μήνυμά σας μεταφέρθηκε στον άλλο παίκτη.} \\
\\ 
\\ 
\texttt{Αυτή τη στιγμή παίζετε τον γύρο 9 από συνολικά 12.} \\
\texttt{Έχετε συσσωρεύσει συλλογικά 4 από 8 ποινές.} \\
\texttt{Ο άλλος παίκτης έστειλε αυτό το μήνυμα:} \\
\\ 
\\ 
\texttt{"Έλαβα ποινή! Στον πίνακά μου το (5,7) έχει '│', δεν είναι άδειο. Λάθος μου που δεν το έλεγξα. Στη σειρά 7 του πίνακά μου: (1,7) άδειο, (2,7) άδειο, (3,7) άδειο, (4,7) C, (5,7) '│', (6,7) άδειο, (7,7) άδειο. Άρα για C, τα άδεια είναι (1,7), (2,7), (3,7), (6,7), (7,7). Εσύ είπες ότι (1,7), (2,7), (6,7) είναι άδεια. Ας διαλέξουμε ένα από αυτά. Προτείνω C στο (6,7). Είναι άδειο στον δικό σου;"} \\
\\ 
\\ 
\texttt{Ποια είναι η επόμενη εντολή σας;} \\
            }
        }
    }
    & \\ \\

    \theutterance \stepcounter{utterance}  
    & & & \multicolumn{4}{p{0.6\linewidth}}{
        \cellcolor[rgb]{0.9,0.9,0.9}{
            \makecell[{{p{\linewidth}}}]{
                \texttt{\tiny{[GM$\langle$P2]}}
                \texttt{ΠΕΣ: Στον πίνακά μου, στη σειρά 7: (1,7) άδειο, (2,7) άδειο, (3,7) '│', (4,7) άδειο, (5,7) άδειο, (6,7) άδειο, (7,7) C. Το (6,7) είναι άδειο! Συμφωνώ για C στο (6,7). Νέοι στόχοι: C(6,7), L(1,4), P(7,5). Θα μετακινήσω τώρα το C στο (6,7).} \\
            }
        }
    }
     \\ \\

    \theutterance \stepcounter{utterance}  
    & & & \multicolumn{2}{p{0.3\linewidth}}{
        \cellcolor[rgb]{0.9,0.9,0.9}{
            \makecell[{{p{\linewidth}}}]{
                \texttt{\tiny{[GM$|$GM]}}
                \texttt{ΠΕΣ: Στον πίνακά μου, στη σειρά 7: (1,7) άδειο, (2,7) άδειο, (3,7) '│', (4,7) άδειο, (5,7) άδειο, (6,7) άδειο, (7,7) C. Το (6,7) είναι άδειο! Συμφωνώ για C στο (6,7). Νέοι στόχοι: C(6,7), L(1,4), P(7,5). Θα μετακινήσω τώρα το C στο (6,7).} \\
            }
        }
    }
    & & \\ \\

    \theutterance \stepcounter{utterance}  
    & & \multicolumn{4}{p{0.6\linewidth}}{
        \cellcolor[rgb]{0.9,0.9,0.9}{
            \makecell[{{p{\linewidth}}}]{
                \texttt{\tiny{[P1$\langle$GM]}}
                \texttt{Το μήνυμά σας μεταφέρθηκε στον άλλο παίκτη.} \\
\\ 
\\ 
\texttt{Αυτή τη στιγμή παίζετε τον γύρο 9 από συνολικά 12.} \\
\texttt{Έχετε συσσωρεύσει συλλογικά 4 από 8 ποινές.} \\
\texttt{Ο άλλος παίκτης έστειλε αυτό το μήνυμα:} \\
\\ 
\\ 
\texttt{"Στον πίνακά μου, στη σειρά 7: (1,7) άδειο, (2,7) άδειο, (3,7) '│', (4,7) άδειο, (5,7) άδειο, (6,7) άδειο, (7,7) C. Το (6,7) είναι άδειο! Συμφωνώ για C στο (6,7). Νέοι στόχοι: C(6,7), L(1,4), P(7,5). Θα μετακινήσω τώρα το C στο (6,7)."} \\
\\ 
\\ 
\texttt{Ποια είναι η επόμενη εντολή σας;} \\
            }
        }
    }
    & \\ \\

    \theutterance \stepcounter{utterance}  
    & \multicolumn{4}{p{0.6\linewidth}}{
        \cellcolor[rgb]{0.9,0.9,0.9}{
            \makecell[{{p{\linewidth}}}]{
                \texttt{\tiny{[P1$\rangle$GM]}}
                \texttt{ΚΙΝΗΣΗ: C, (6,7)} \\
            }
        }
    }
    & & \\ \\

    \theutterance \stepcounter{utterance}  
    & & & \multicolumn{2}{p{0.3\linewidth}}{
        \cellcolor[rgb]{0.9,0.9,0.9}{
            \makecell[{{p{\linewidth}}}]{
                \texttt{\tiny{[GM$|$GM]}}
                \texttt{ΚΙΝΗΣΗ: C, (6,7)} \\
            }
        }
    }
    & & \\ \\

    \theutterance \stepcounter{utterance}  
    & & & \multicolumn{2}{p{0.3\linewidth}}{
        \cellcolor[rgb]{0.9,0.9,0.9}{
            \makecell[{{p{\linewidth}}}]{
                \texttt{\tiny{[GM$|$GM]}}
                \texttt{C moved to (6, 7)} \\
            }
        }
    }
    & & \\ \\

    \theutterance \stepcounter{utterance}  
    & & \multicolumn{4}{p{0.6\linewidth}}{
        \cellcolor[rgb]{0.9,0.9,0.9}{
            \makecell[{{p{\linewidth}}}]{
                \texttt{\tiny{[GM$\rangle$P2]}}
                \texttt{Το μήνυμά σας μεταφέρθηκε στον άλλο παίκτη.} \\
\\ 
\\ 
\texttt{Αυτή τη στιγμή παίζετε τον γύρο 10 από συνολικά 12.} \\
\texttt{Έχετε συσσωρεύσει συλλογικά 4 από 8 ποινές.} \\
\texttt{Ο άλλος παίκτης μετακίνησε ένα αντικείμενο στο πλέγμα του.} \\
\\ 
\\ 
\texttt{Ποια είναι η επόμενη εντολή σας;} \\
            }
        }
    }
    & \\ \\

    \theutterance \stepcounter{utterance}  
    & & & \multicolumn{4}{p{0.6\linewidth}}{
        \cellcolor[rgb]{0.9,0.9,0.9}{
            \makecell[{{p{\linewidth}}}]{
                \texttt{\tiny{[GM$\langle$P2]}}
                \texttt{ΚΙΝΗΣΗ: C, (6,7)} \\
            }
        }
    }
     \\ \\

    \theutterance \stepcounter{utterance}  
    & & & \multicolumn{2}{p{0.3\linewidth}}{
        \cellcolor[rgb]{0.9,0.9,0.9}{
            \makecell[{{p{\linewidth}}}]{
                \texttt{\tiny{[GM$|$GM]}}
                \texttt{ΚΙΝΗΣΗ: C, (6,7)} \\
            }
        }
    }
    & & \\ \\

    \theutterance \stepcounter{utterance}  
    & & & \multicolumn{2}{p{0.3\linewidth}}{
        \cellcolor[rgb]{0.9,0.9,0.9}{
            \makecell[{{p{\linewidth}}}]{
                \texttt{\tiny{[GM$|$GM]}}
                \texttt{C moved to (6, 7)} \\
            }
        }
    }
    & & \\ \\

    \theutterance \stepcounter{utterance}  
    & & \multicolumn{4}{p{0.6\linewidth}}{
        \cellcolor[rgb]{0.9,0.9,0.9}{
            \makecell[{{p{\linewidth}}}]{
                \texttt{\tiny{[P1$\langle$GM]}}
                \texttt{Μεταφέρθηκε επιτυχώς το 'C' στο (6,7). Το ενημερωμένο πλέγμα σας μοιάζει ως εξής:} \\
\texttt{```} \\
\texttt{ 1234567} \\
\texttt{╔══╤═╤══╗} \\
\texttt{║◌◌│◌│◌◌║ 1} \\
\texttt{║◌◌└─┼──╢ 2} \\
\texttt{║◌◌◌◌│◌◌║ 3} \\
\texttt{║◌┌─┬┴┬─╢ 4} \\
\texttt{╟─┤◌│◌│P║ 5} \\
\texttt{║◌└┬┴┬┘L║ 6} \\
\texttt{║◌◌│◌│C◌║ 7} \\
\texttt{╚══╧═╧══╝} \\
\texttt{```} \\
\texttt{Αυτή τη στιγμή παίζετε τον γύρο 10 από συνολικά 12.} \\
\texttt{Έχετε συσσωρεύσει συλλογικά 4 από 8 ποινές.} \\
\texttt{Ο άλλος παίκτης μετακίνησε ένα αντικείμενο στο πλέγμα του.} \\
\\ 
\\ 
\texttt{Ποια είναι η επόμενη εντολή σας;} \\
            }
        }
    }
    & \\ \\

    \theutterance \stepcounter{utterance}  
    & \multicolumn{4}{p{0.6\linewidth}}{
        \cellcolor[rgb]{0.9,0.9,0.9}{
            \makecell[{{p{\linewidth}}}]{
                \texttt{\tiny{[P1$\rangle$GM]}}
                \texttt{ΠΕΣ: Το C μου είναι τώρα στο (6,7)! Απομένει το L. Στόχος L στο (1,4). Στον πίνακά μου το (1,4) είναι άδειο. Εσύ επιβεβαίωσες ότι είναι άδειο. Θα μετακινήσω το L από το (7,6) στο (1,4). Εσύ έτοιμος για το L;} \\
            }
        }
    }
    & & \\ \\

    \theutterance \stepcounter{utterance}  
    & & & \multicolumn{2}{p{0.3\linewidth}}{
        \cellcolor[rgb]{0.9,0.9,0.9}{
            \makecell[{{p{\linewidth}}}]{
                \texttt{\tiny{[GM$|$GM]}}
                \texttt{ΠΕΣ: Το C μου είναι τώρα στο (6,7)! Απομένει το L. Στόχος L στο (1,4). Στον πίνακά μου το (1,4) είναι άδειο. Εσύ επιβεβαίωσες ότι είναι άδειο. Θα μετακινήσω το L από το (7,6) στο (1,4). Εσύ έτοιμος για το L;} \\
            }
        }
    }
    & & \\ \\

    \theutterance \stepcounter{utterance}  
    & & \multicolumn{4}{p{0.6\linewidth}}{
        \cellcolor[rgb]{0.9,0.9,0.9}{
            \makecell[{{p{\linewidth}}}]{
                \texttt{\tiny{[GM$\rangle$P2]}}
                \texttt{Μεταφέρθηκε επιτυχώς το 'C' στο (6,7). Το ενημερωμένο πλέγμα σας μοιάζει ως εξής:} \\
\texttt{```} \\
\texttt{ 1234567} \\
\texttt{╔══╤═╤══╗} \\
\texttt{║◌◌│◌│◌◌║ 1} \\
\texttt{║◌L└─┼──╢ 2} \\
\texttt{║◌◌◌◌│◌◌║ 3} \\
\texttt{║◌┌─┬┴┬─╢ 4} \\
\texttt{╟─┤◌│◌│P║ 5} \\
\texttt{║◌└┬┴┬┘◌║ 6} \\
\texttt{║◌◌│◌│C◌║ 7} \\
\texttt{╚══╧═╧══╝} \\
\texttt{```} \\
\texttt{Αυτή τη στιγμή παίζετε τον γύρο 11 από συνολικά 12.} \\
\texttt{Έχετε συσσωρεύσει συλλογικά 4 από 8 ποινές.} \\
\texttt{Ο άλλος παίκτης έστειλε αυτό το μήνυμα:} \\
\\ 
\\ 
\texttt{"Το C μου είναι τώρα στο (6,7)! Απομένει το L. Στόχος L στο (1,4). Στον πίνακά μου το (1,4) είναι άδειο. Εσύ επιβεβαίωσες ότι είναι άδειο. Θα μετακινήσω το L από το (7,6) στο (1,4). Εσύ έτοιμος για το L;"} \\
\\ 
\\ 
\texttt{Ποια είναι η επόμενη εντολή σας;} \\
            }
        }
    }
    & \\ \\

    \theutterance \stepcounter{utterance}  
    & & & \multicolumn{4}{p{0.6\linewidth}}{
        \cellcolor[rgb]{0.9,0.9,0.9}{
            \makecell[{{p{\linewidth}}}]{
                \texttt{\tiny{[GM$\langle$P2]}}
                \texttt{ΚΙΝΗΣΗ: L, (1,4)} \\
            }
        }
    }
     \\ \\

    \theutterance \stepcounter{utterance}  
    & & & \multicolumn{2}{p{0.3\linewidth}}{
        \cellcolor[rgb]{0.9,0.9,0.9}{
            \makecell[{{p{\linewidth}}}]{
                \texttt{\tiny{[GM$|$GM]}}
                \texttt{ΚΙΝΗΣΗ: L, (1,4)} \\
            }
        }
    }
    & & \\ \\

    \theutterance \stepcounter{utterance}  
    & & & \multicolumn{2}{p{0.3\linewidth}}{
        \cellcolor[rgb]{0.9,0.9,0.9}{
            \makecell[{{p{\linewidth}}}]{
                \texttt{\tiny{[GM$|$GM]}}
                \texttt{L moved to (1, 4)} \\
            }
        }
    }
    & & \\ \\

    \theutterance \stepcounter{utterance}  
    & & \multicolumn{4}{p{0.6\linewidth}}{
        \cellcolor[rgb]{0.9,0.9,0.9}{
            \makecell[{{p{\linewidth}}}]{
                \texttt{\tiny{[P1$\langle$GM]}}
                \texttt{Το μήνυμά σας μεταφέρθηκε στον άλλο παίκτη.} \\
\\ 
\\ 
\texttt{Αυτή τη στιγμή παίζετε τον γύρο 11 από συνολικά 12.} \\
\texttt{Έχετε συσσωρεύσει συλλογικά 4 από 8 ποινές.} \\
\texttt{Ο άλλος παίκτης μετακίνησε ένα αντικείμενο στο πλέγμα του.} \\
\\ 
\\ 
\texttt{Ποια είναι η επόμενη εντολή σας;} \\
            }
        }
    }
    & \\ \\

    \theutterance \stepcounter{utterance}  
    & \multicolumn{4}{p{0.6\linewidth}}{
        \cellcolor[rgb]{0.9,0.9,0.9}{
            \makecell[{{p{\linewidth}}}]{
                \texttt{\tiny{[P1$\rangle$GM]}}
                \texttt{ΚΙΝΗΣΗ: L, (1,4)} \\
            }
        }
    }
    & & \\ \\

    \theutterance \stepcounter{utterance}  
    & & & \multicolumn{2}{p{0.3\linewidth}}{
        \cellcolor[rgb]{0.9,0.9,0.9}{
            \makecell[{{p{\linewidth}}}]{
                \texttt{\tiny{[GM$|$GM]}}
                \texttt{ΚΙΝΗΣΗ: L, (1,4)} \\
            }
        }
    }
    & & \\ \\

    \theutterance \stepcounter{utterance}  
    & & & \multicolumn{2}{p{0.3\linewidth}}{
        \cellcolor[rgb]{0.9,0.9,0.9}{
            \makecell[{{p{\linewidth}}}]{
                \texttt{\tiny{[GM$|$GM]}}
                \texttt{L moved to (1, 4)} \\
            }
        }
    }
    & & \\ \\

    \theutterance \stepcounter{utterance}  
    & & \multicolumn{4}{p{0.6\linewidth}}{
        \cellcolor[rgb]{0.9,0.9,0.9}{
            \makecell[{{p{\linewidth}}}]{
                \texttt{\tiny{[GM$\rangle$P2]}}
                \texttt{Μεταφέρθηκε επιτυχώς το 'L' στο (1,4). Το ενημερωμένο πλέγμα σας μοιάζει ως εξής:} \\
\texttt{```} \\
\texttt{ 1234567} \\
\texttt{╔══╤═╤══╗} \\
\texttt{║◌◌│◌│◌◌║ 1} \\
\texttt{║◌◌└─┼──╢ 2} \\
\texttt{║◌◌◌◌│◌◌║ 3} \\
\texttt{║L┌─┬┴┬─╢ 4} \\
\texttt{╟─┤◌│◌│P║ 5} \\
\texttt{║◌└┬┴┬┘◌║ 6} \\
\texttt{║◌◌│◌│C◌║ 7} \\
\texttt{╚══╧═╧══╝} \\
\texttt{```} \\
\texttt{Αυτή τη στιγμή παίζετε τον γύρο 12 από συνολικά 12.} \\
\texttt{Έχετε συσσωρεύσει συλλογικά 4 από 8 ποινές.} \\
\texttt{Ο άλλος παίκτης μετακίνησε ένα αντικείμενο στο πλέγμα του.} \\
\\ 
\\ 
\texttt{Ποια είναι η επόμενη εντολή σας;} \\
            }
        }
    }
    & \\ \\

    \theutterance \stepcounter{utterance}  
    & & & \multicolumn{4}{p{0.6\linewidth}}{
        \cellcolor[rgb]{0.9,0.9,0.9}{
            \makecell[{{p{\linewidth}}}]{
                \texttt{\tiny{[GM$\langle$P2]}}
                \texttt{ΠΕΣ: Το L μου είναι τώρα στο (1,4)! Ο πίνακάς μου έχει: P στο (7,5), C στο (6,7), L στο (1,4). Εσύ έχεις όλα τα αντικείμενα στις θέσεις{-}στόχους; Αν ναι, τελείωσες;} \\
            }
        }
    }
     \\ \\

    \theutterance \stepcounter{utterance}  
    & & & \multicolumn{2}{p{0.3\linewidth}}{
        \cellcolor[rgb]{0.9,0.9,0.9}{
            \makecell[{{p{\linewidth}}}]{
                \texttt{\tiny{[GM$|$GM]}}
                \texttt{ΠΕΣ: Το L μου είναι τώρα στο (1,4)! Ο πίνακάς μου έχει: P στο (7,5), C στο (6,7), L στο (1,4). Εσύ έχεις όλα τα αντικείμενα στις θέσεις{-}στόχους; Αν ναι, τελείωσες;} \\
            }
        }
    }
    & & \\ \\

    \theutterance \stepcounter{utterance}  
    & & & \multicolumn{2}{p{0.3\linewidth}}{
        \cellcolor[rgb]{0.9,0.9,0.9}{
            \makecell[{{p{\linewidth}}}]{
                \texttt{\tiny{[GM$|$GM]}}
                \texttt{Maximum number of rounds reached} \\
            }
        }
    }
    & & \\ \\

    \theutterance \stepcounter{utterance}  
    & & & \multicolumn{2}{p{0.3\linewidth}}{
        \cellcolor[rgb]{0.9,0.9,0.9}{
            \makecell[{{p{\linewidth}}}]{
                \texttt{\tiny{[GM$|$GM]}}
                \texttt{Initial states:} \\
\texttt{```} \\
\texttt{Player 1:} \\
\texttt{ 1234567} \\
\texttt{╔══╤═╤══╗} \\
\texttt{║◌◌│◌│◌◌║ 1} \\
\texttt{║◌◌└─┼──╢ 2} \\
\texttt{║◌◌◌◌│P◌║ 3} \\
\texttt{║◌┌─┬┴┬─╢ 4} \\
\texttt{╟─┤◌│◌│◌║ 5} \\
\texttt{║◌└┬┴┬┘L║ 6} \\
\texttt{║◌◌│C│◌◌║ 7} \\
\texttt{╚══╧═╧══╝} \\
\\ 
\\ 
\texttt{Player 2:} \\
\texttt{ 1234567} \\
\texttt{╔══╤═╤══╗} \\
\texttt{║◌P│◌│◌◌║ 1} \\
\texttt{║◌L└─┼──╢ 2} \\
\texttt{║◌◌◌◌│◌◌║ 3} \\
\texttt{║◌┌─┬┴┬─╢ 4} \\
\texttt{╟─┤◌│◌│◌║ 5} \\
\texttt{║◌└┬┴┬┘◌║ 6} \\
\texttt{║◌◌│◌│◌C║ 7} \\
\texttt{╚══╧═╧══╝} \\
\\ 
\texttt{```} \\
            }
        }
    }
    & & \\ \\

    \theutterance \stepcounter{utterance}  
    & & & \multicolumn{2}{p{0.3\linewidth}}{
        \cellcolor[rgb]{0.9,0.9,0.9}{
            \makecell[{{p{\linewidth}}}]{
                \texttt{\tiny{[GM$|$GM]}}
                \texttt{End states:} \\
\texttt{```} \\
\texttt{Player 1:} \\
\texttt{ 1234567} \\
\texttt{╔══╤═╤══╗} \\
\texttt{║◌◌│◌│◌◌║ 1} \\
\texttt{║◌◌└─┼──╢ 2} \\
\texttt{║◌◌◌◌│◌◌║ 3} \\
\texttt{║L┌─┬┴┬─╢ 4} \\
\texttt{╟─┤◌│◌│P║ 5} \\
\texttt{║◌└┬┴┬┘◌║ 6} \\
\texttt{║◌◌│◌│C◌║ 7} \\
\texttt{╚══╧═╧══╝} \\
\\ 
\\ 
\texttt{Player 2:} \\
\texttt{ 1234567} \\
\texttt{╔══╤═╤══╗} \\
\texttt{║◌◌│◌│◌◌║ 1} \\
\texttt{║◌◌└─┼──╢ 2} \\
\texttt{║◌◌◌◌│◌◌║ 3} \\
\texttt{║L┌─┬┴┬─╢ 4} \\
\texttt{╟─┤◌│◌│P║ 5} \\
\texttt{║◌└┬┴┬┘◌║ 6} \\
\texttt{║◌◌│◌│C◌║ 7} \\
\texttt{╚══╧═╧══╝} \\
\\ 
\texttt{```} \\
            }
        }
    }
    & & \\ \\

    \theutterance \stepcounter{utterance}  
    & & & \multicolumn{2}{p{0.3\linewidth}}{
        \cellcolor[rgb]{0.9,0.9,0.9}{
            \makecell[{{p{\linewidth}}}]{
                \texttt{\tiny{[GM$|$GM]}}
                \texttt{* success: True} \\
\texttt{* lose: False} \\
\texttt{* aborted: False} \\
\texttt{{-}{-}{-}{-}{-}{-}{-}} \\
\texttt{* Shifts: 2.00} \\
\texttt{* Max Shifts: 4.00} \\
\texttt{* Min Shifts: 2.00} \\
\texttt{* End Distance Sum: 0.00} \\
\texttt{* Init Distance Sum: 13.88} \\
\texttt{* Expected Distance Sum: 12.57} \\
\texttt{* Penalties: 4.00} \\
\texttt{* Max Penalties: 8.00} \\
\texttt{* Rounds: 12.00} \\
\texttt{* Max Rounds: 12.00} \\
\texttt{* Object Count: 3.00} \\
            }
        }
    }
    & & \\ \\

\end{supertabular}
}

\end{document}
